\documentclass[11pt,letterpaper]{article}

% ============================================================
% PAQUETES
% ============================================================
\usepackage[utf8]{inputenc}
\usepackage[T1]{fontenc}
\usepackage[spanish,es-tabla,es-nodecimaldot]{babel}
\usepackage{mathptmx} % Tipografía formal estilo Times
\usepackage[left=2.2cm, right=2.2cm, top=2.5cm, bottom=2.5cm, headheight=28pt]{geometry}
\usepackage{xcolor}
\usepackage{graphicx}
\usepackage{tikz}
\usepackage{eso-pic}
\usepackage{fancyhdr}
\usepackage{booktabs}
\usepackage{longtable}
\usepackage{tabularx}
\usepackage{array}
\usepackage{enumitem}
\usepackage{tcolorbox}
\usepackage{amssymb}
\usepackage{pifont}
\usepackage{setspace}
\usepackage{parskip}
\usepackage{hyperref}

\tcbuselibrary{skins,breakable}

% ============================================================
% PALETA DE COLORES (Inspirada en Mezcal Mil Sensaciones)
% ============================================================
\definecolor{mezcalDark}{HTML}{1E1B18}     % Carbón orgánico
\definecolor{mezcalEarth}{HTML}{5C3A21}    % Tierra / Madera
\definecolor{mezcalAgave}{HTML}{2E5A44}    % Verde Agave
\definecolor{mezcalGold}{HTML}{B8860B}     % Ámbar / Dorado
\definecolor{mezcalCream}{HTML}{FAF7F2}    % Fondo suave crema
\definecolor{alertBg}{HTML}{FFF8E7}        % Fondo para pendientes
\definecolor{watermarkColor}{HTML}{D0C8BE} % Gris/Tierra muy tenue para marca de agua

\hypersetup{
    colorlinks=true,
    linkcolor=mezcalEarth,
    urlcolor=mezcalAgave,
    citecolor=mezcalEarth,
    pdftitle={SRS y Cuestionario de Levantamiento - Mezcal Mil Sensaciones},
    pdfauthor={Ing. en Computacion Jose Ramon Aragon Toledo (Clave DGP 521344)}
}

% ============================================================
% MARCA DE AGUA CON DATOS DEL AUTOR EN CADA PÁGINA
% ============================================================
\AddToShipoutPictureBG{%
    \begin{tikzpicture}[remember picture, overlay]
        \node[
            rotate=38,
            scale=1.45,
            text=watermarkColor,
            opacity=0.28,
            align=center,
            font=\bfseries\sffamily
        ] at (current page.center) {
            ING. EN COMPUTACIÓN JOSE RAMON ARAGON TOLEDO\\[0.25em]
            {\large Ingeniero en Computación \quad|\quad Clave DGP: 521344}\\[0.25em]
            {\normalsize aatr010423@gmail.com}\\[0.2em]
            {\normalsize Desarrollo y Mantenimiento: Ing. en Computación Jose Ramon Aragon Toledo}
        };
    \end{tikzpicture}%
}

% ============================================================
% ENCABEZADOS Y PIES DE PÁGINA
% ============================================================
\pagestyle{fancy}
\fancyhf{}
\fancyhead[L]{\small\textbf{\textcolor{mezcalEarth}{SRS — Mezcal Artesanal ``Mil Sensaciones''}}}
\fancyhead[R]{\small\textcolor{mezcalAgave}{IEEE 830 / ISO 29148 \quad|\quad v1.2}}
\fancyfoot[L]{\footnotesize\textcolor{gray}{Autor: Ing. en Computación Jose Ramon Aragon Toledo (Clave DGP: 521344) — \texttt{aatr010423@gmail.com}}}
\fancyfoot[R]{\footnotesize\textbf{Página \thepage}}
\renewcommand{\headrulewidth}{0.5pt}
\renewcommand{\footrulewidth}{0.5pt}

% ============================================================
% COMANDOS AUXILIARES PARA CUESTIONARIO Y ESPACIOS RELLENABLES
% ============================================================
\newcommand{\cmark}{\textcolor{mezcalAgave}{\ding{52}}}
\newcommand{\uchk}{$\square$}
\newcommand{\chk}{\textcolor{mezcalAgave}{$\boxtimes$}}
\newcommand{\uradio}{$\bigcirc$}
\newcommand{\radio}{\textcolor{mezcalAgave}{\textbf{$\odot$}}}

% Línea para rellenar a mano o digitalmente
\newcommand{\fillline}[1][1.0\linewidth]{%
    \par\vspace{0.35em}\noindent\textcolor{gray!70}{\rule{#1}{0.5pt}}\par\vspace{0.25em}%
}

% Caja de pregunta pendiente con espacio para respuesta
\newtcolorbox{pendingbox}[1]{
    enhanced,
    breakable,
    colback=alertBg,
    colframe=mezcalGold,
    coltitle=white,
    fonttitle=\bfseries,
    title={#1},
    arc=3pt,
    boxrule=0.8pt,
    left=8pt, right=8pt, top=6pt, bottom=6pt
}

% Caja de respuesta registrada del cuestionario
\newtcolorbox{answerbox}{
    enhanced,
    breakable,
    colback=mezcalCream,
    colframe=mezcalAgave!60,
    arc=2pt,
    boxrule=0.5pt,
    left=8pt, right=8pt, top=5pt, bottom=5pt
}

\begin{document}

% ============================================================
% PORTADA
% ============================================================
\begin{titlepage}
    \begin{flushright}
        \small
        \textbf{Ingeniería de Software y Desarrollo Web Profesional}\\
        \textbf{Autor:} Ing. en Computación Jose Ramon Aragon Toledo\\
        \textbf{Título:} Ingeniero en Computación (Clave DGP: \texttt{521344})\\
        \textbf{Correo:} \href{mailto:aatr010423@gmail.com}{aatr010423@gmail.com}\\
        \textbf{Equipo de Desarrollo:} Ing. en Computación Jose Ramon Aragon Toledo
    \end{flushright}

    \vspace{1.2cm}

    \begin{center}
        {\color{mezcalEarth}\rule{\linewidth}{2pt}}\\[0.4cm]
        {\Huge\bfseries\textcolor{mezcalDark}{Especificación de Requerimientos de Software (SRS)}}\\[0.3cm]
        {\LARGE\bfseries\textcolor{mezcalEarth}{y Levantamiento de Información del Cliente}}\\[0.3cm]
        {\Large\textcolor{mezcalAgave}{Plataforma Web Oficial — Mezcal Artesanal ``Mil Sensaciones''}}\\[0.2cm]
        {\color{mezcalEarth}\rule{\linewidth}{2pt}}
    \end{center}

    \vspace{0.8cm}

    \begin{tcolorbox}[colback=mezcalCream, colframe=mezcalEarth, title=\textbf{Ficha de Control del Documento}, fonttitle=\bfseries\large]
        \renewcommand{\arraystretch}{1.35}
        \begin{tabularx}{\linewidth}{@{} l >{\raggedright\arraybackslash}X @{}}
            \textbf{Estándar de Referencia:} & IEEE 830 / ISO/IEC/IEEE 29148 \\
            \textbf{Versión del Documento:} & 1.2 (Integra SRS + Cuestionario de Levantamiento + Formato de Pendientes) \\
            \textbf{Fecha de Emisión:} & 23 de octubre de 2026 (revisión 1.2: 26 de octubre de 2026) \\
            \textbf{Documento Relacionado:} & Plan de Desarrollo Técnico y Presupuesto: \$7,550.00 MXN a 8 semanas \\
            \textbf{Cliente / Marca:} & Mezcal Mil Sensaciones \\
            \textbf{Origen / Palenque:} & El Gramal, Nejapa de Madero, Yautepec, Oaxaca, México \\
            \textbf{Maestro Mezcalero:} & Damián Olivera Díaz (2.ª Generación; marca en 3.ª Generación familiar) \\
            \textbf{Responsable / Autor:} & Ing. en Computación Jose Ramon Aragon Toledo (Clave DGP: \texttt{521344}) \\
            \textbf{Equipo de Desarrollo:} & Ing. en Computación Jose Ramon Aragon Toledo (\texttt{aatr010423@gmail.com}) \\
        \end{tabularx}
    \end{tcolorbox}

    \vfill

    \begin{center}
        \footnotesize\textcolor{gray}{Documento técnico confidencial con marca de agua de autoría registrada a nombre de José Ramón Aragón Toledo.\\ Incluye secciones interactivas/imprimibles para completar las preguntas y activos pendientes con el cliente.}
    \end{center}
\end{titlepage}

% ============================================================
% TABLA DE CONTENIDO
% ============================================================
\tableofcontents
\newpage

% ============================================================
% HISTORIAL DE CAMBIOS
% ============================================================
\section*{Historial de Cambios del Documento}
\addcontentsline{toc}{section}{Historial de Cambios del Documento}

\renewcommand{\arraystretch}{1.35}
\begin{tabularx}{\linewidth}{@{} l l >{\raggedright\arraybackslash}X @{}}
\toprule
\textbf{Versión} & \textbf{Fecha} & \textbf{Descripción del cambio} \\
\midrule
1.0 & 23/10/2026 & Emisión inicial: SRS del sitio institucional y registro del cuestionario de levantamiento. \\
\midrule
1.1 & 26/10/2026 & Alineación con el \textit{Plan de Desarrollo Técnico y Presupuesto} (\$6,050.00 MXN a 6 semanas, versión vigente en esa fecha). Se agregan a la Fase 1 los módulos comerciales contratados (RF-16 a RF-25), la arquitectura técnica con backend en Go, la matriz de roles y permisos, los criterios de aceptación, el alcance excluido, los supuestos, nuevos requerimientos no funcionales y la actualización del estado real del prototipo. \\
\midrule
1.2 & 26/10/2026 & Se incorporan a la Fase 1 los pagos en línea (RF-26) y la facturación electrónica CFDI 4.0 (RF-27), que antes eran Fase 2. El plazo pasa a 8 semanas y el presupuesto a \$7,550.00 MXN (\$1,550.00 de infraestructura y \$6,000.00 de honorarios en 4 quincenas). Se actualizan roles, arquitectura, supuestos fiscales, exclusiones y formularios de pendientes. \\
\bottomrule
\end{tabularx}

% ============================================================
% SECCIÓN 1: INTRODUCCIÓN
% ============================================================
\section{Introducción}

\subsection{Propósito del Documento}
El presente documento de \textbf{Especificación de Requerimientos de Software (SRS)} define de manera formal, completa y verificable los requerimientos funcionales, no funcionales, de contenido, diseño, legales y operativos para el desarrollo y lanzamiento de la plataforma web oficial de \textbf{Mezcal Mil Sensaciones}. Junto con el \textit{Plan de Desarrollo Técnico y Presupuesto}, es la referencia técnica para aceptar cada entrega. Adicionalmente, integra el registro íntegro del cuestionario de levantamiento aplicado al cliente y proporciona los \textbf{espacios y formularios estructurados} para responder las preguntas que quedaron sin contestar y recopilar los activos pendientes.

\subsection{Alcance del Producto (Fase 1 vs. Fase 2)}

\textbf{Fase 1 — Lanzamiento contratado (8 semanas, incluido en el presupuesto de \$7,550.00 MXN):}
\begin{itemize}[leftmargin=1.5em]
    \item \textbf{Bloque A. Sitio institucional y de marca:} sitio interactivo, bilingüe (Español / Inglés) y responsivo con la historia familiar de tres generaciones, el perfil del Maestro Mezcalero \textbf{Damián Olivera Díaz}, el proceso artesanal, las fichas técnicas de cada expresión con descarga en PDF, galería multimedia, ubicación en Google Maps, canales de contacto (WhatsApp, formulario y redes sociales), verificación de mayoría de edad (+18) y apartados legales (RF-01 a RF-15).
    \item \textbf{Bloque B. Plataforma comercial y de gestión:} carrito de compras, registro de pedidos, pagos en línea con tarjeta, OXXO y SPEI, facturación electrónica CFDI 4.0, cuentas de usuario con tres roles, dashboard administrativo, terminal de punto de venta (POS), seguimiento de pedidos, chatbot de WhatsApp, chat en vivo dentro del sitio y correos transaccionales (RF-16 a RF-27).
\end{itemize}

\textbf{Fase 2 — Evolutivos futuros (fuera del presupuesto actual, se cotizan por separado):}
\begin{itemize}[leftmargin=1.5em]
    \item Integración automática con paqueterías (cálculo de envío y generación de guías por API).
    \item Sistema de reservas para catas o tours al palenque en El Gramal.
    \item Portal B2B para distribuidores e importadores con precios especiales.
    \item Gestor de contenido editorial (CMS) para que el cliente edite textos e imágenes del sitio sin programador.
    \item Blog cultural del mezcal e integración con publicaciones recientes de Instagram.
    \item Integración con marketplaces (Amazon, MercadoLibre) y aplicación móvil nativa.
\end{itemize}

\subsection{Exclusiones Explícitas de la Fase 1}
Para evitar ambigüedades, los siguientes puntos \textbf{no} forman parte de la Fase 1:
\begin{itemize}[leftmargin=1.5em]
    \item \textbf{Asesoría contable y fiscal.} El equipo implementa la facturación con la configuración fiscal (régimen, IVA, IEPS y claves del SAT) que defina el contador del cliente. No incluye declaraciones, contabilidad ni trámites ante el SAT.
    \item \textbf{Comisiones de pago y timbrado.} Las comisiones de la pasarela por cada venta y el costo de cada factura timbrada corren a cargo del cliente.
    \item \textbf{Venta directa a consumidores en Estados Unidos.} El sitio informa y capta contactos de importadores. La venta en EE.~UU. requiere un importador autorizado en ese país.
    \item \textbf{Modo sin conexión del POS.} La terminal POS requiere conexión a internet. No se incluyen impresoras de tickets, lectores de código ni terminales bancarias.
    \item \textbf{Inteligencia artificial conversacional.} El chatbot de WhatsApp funciona con menús y palabras clave predefinidas.
    \item \textbf{Producción fotográfica o de video profesional.} El cliente entrega el material. El equipo de desarrollo solo genera los videos \textit{loopable} de botellas con IA a partir de ese material.
    \item \textbf{Validación jurídica de los textos legales.} El equipo redacta el Aviso de Privacidad y los Términos con base en la ley vigente. Se recomienda que el cliente los haga revisar por un abogado.
    \item \textbf{Costos recurrentes posteriores al mes 12.} El servidor, la renovación del dominio y el consumo de mensajes de WhatsApp que rebase el límite gratuito quedan a cargo del cliente, según el Plan de Desarrollo.
\end{itemize}

\subsection{Público Objetivo y Mercados}
\begin{enumerate}[leftmargin=1.5em]
    \item \textbf{Consumidor final (+18 años):} Busca mezcal artesanal auténtico con identidad, carácter e historia familiar.
    \item \textbf{Bares, restaurantes y centros de coctelería (On-Trade):} Buscan fichas técnicas, graduación alcohólica, volumen por lote, perfil sensorial y contacto comercial directo.
    \item \textbf{Importadores internacionales (EE. UU.):} Requieren información regulatoria clara (Código NOM-070 / CRM, registro IMPI, capacidad por lote) y navegación 100\% en inglés.
    \item \textbf{Turistas y visitantes en Oaxaca / México:} Buscan conocer el origen, el palenque y dónde adquirir las botellas.
\end{enumerate}
\textbf{Mercados prioritarios:} Oaxaca (México) y Estados Unidos (exportación/importación).

\subsection{Definiciones, Acrónimos y Abreviaturas}
\renewcommand{\arraystretch}{1.25}
\begin{longtable}{@{} >{\raggedright\arraybackslash}p{3.4cm} >{\raggedright\arraybackslash}p{12.4cm} @{}}
\toprule
\textbf{Término} & \textbf{Significado} \\
\midrule
\endhead
SRS & Especificación de Requerimientos de Software (\textit{Software Requirements Specification}). \\
RF / RNF & Requerimiento Funcional / Requerimiento No Funcional. \\
MVP & Producto mínimo viable: primera versión publicable del sistema. \\
NOM-070-SCFI-2016 & Norma Oficial Mexicana de bebidas alcohólicas, mezcal: especificaciones y etiquetado. \\
CRM & Consejo Regulador del Mezcal: organismo que certifica el mezcal y emite los marbetes. \\
Marbete & Etiqueta de control fiscal y de certificación adherida a cada botella. \\
IMPI & Instituto Mexicano de la Propiedad Industrial (registro de marca). \\
LFPDPPP & Ley Federal de Protección de Datos Personales en Posesión de los Particulares. \\
LFPC & Ley Federal de Protección al Consumidor (aplica a las ventas a distancia). \\
Derechos ARCO & Derechos de Acceso, Rectificación, Cancelación y Oposición sobre datos personales. \\
POS & Punto de Venta (\textit{Point of Sale}): caja para ventas presenciales. \\
Lote & Producción de una misma destilación, identificada con número, fecha y cantidad de botellas. \\
Folio & Identificador único de un pedido, por ejemplo \texttt{MS-2026-0001}. \\
Go & Lenguaje de programación del backend; se compila como un solo binario ejecutable. \\
API & Interfaz de programación por la cual el sitio se comunica con el servidor. \\
JWT & Token firmado que identifica la sesión del usuario. \\
Argon2id & Algoritmo de cifrado unidireccional para almacenar contraseñas. \\
WebSocket & Canal de comunicación en tiempo real usado por el chat en vivo. \\
VPS & Servidor privado virtual en la nube. \\
Pasarela de pago & Servicio que cobra con tarjeta, OXXO o SPEI y deposita el dinero al negocio (Mercado Pago o Stripe). \\
Webhook & Aviso automático que un servicio externo envía al servidor, por ejemplo cuando se confirma un pago. \\
CFDI 4.0 & Comprobante Fiscal Digital por Internet: la factura electrónica vigente del SAT. \\
PAC & Proveedor Autorizado de Certificación: timbra las facturas ante el SAT. \\
CSD & Certificado de Sello Digital del SAT con el que el negocio firma sus facturas. Es distinto de la e.firma. \\
IEPS & Impuesto Especial sobre Producción y Servicios que grava las bebidas alcohólicas. \\
CDN & Red de distribución de contenido que acelera la entrega de archivos. \\
SPF / DKIM / DMARC & Registros DNS que autentican el correo del dominio y evitan que llegue a spam. \\
LCP / INP / CLS & Métricas de rendimiento web de Google: carga del elemento principal, respuesta a la interacción y estabilidad visual. \\
Video \textit{loopable} & Video cuyo primer y último fotograma coinciden para repetirse sin cortes. \\
\bottomrule
\end{longtable}

\subsection{Documentos de Referencia}
\begin{itemize}[leftmargin=1.5em]
    \item Cuestionario de levantamiento aplicado al cliente (\texttt{Cuestionario\_Mezcal\_Fercho 2.docx}).
    \item Plan de Desarrollo Técnico y Esquema Comercial de Implementación en Producción, actualizado el 26 de octubre de 2026.
    \item IEEE 830-1998 e ISO/IEC/IEEE 29148:2018, ingeniería de requerimientos.
    \item NOM-070-SCFI-2016 (mezcal) y NOM-142-SSA1/SCFI-2014 (etiquetado sanitario de bebidas alcohólicas).
    \item LFPDPPP y LFPC vigentes.
\end{itemize}

\subsection{Supuestos y Dependencias}
\begin{itemize}[leftmargin=1.5em]
    \item \textbf{Entrega de contenidos:} el cliente entrega la información de la Sección~6 a más tardar al final de la semana 2. Si se retrasa, los módulos se publican con contenido provisional marcado como tal y el retraso no se imputa al equipo de desarrollo.
    \item \textbf{Número de WhatsApp:} el cliente proporciona un número telefónico dedicado para la WhatsApp Business Platform. Ese número no puede estar activo al mismo tiempo en la aplicación de WhatsApp del teléfono.
    \item \textbf{Verificación en Meta:} el titular del negocio completa la verificación de su empresa en Meta Business para habilitar el chatbot.
    \item \textbf{Titularidad:} el dominio, las cuentas de los servicios y los datos quedan a nombre del cliente. El equipo de desarrollo conserva acceso técnico durante el proyecto.
    \item \textbf{Cuenta de la pasarela:} el cliente abre su cuenta de Mercado Pago (o Stripe) a nombre del negocio y obtiene la aprobación para vender bebidas alcohólicas antes de la semana 7.
    \item \textbf{Datos fiscales:} antes de la semana 7 el cliente entrega su RFC activo, régimen fiscal, código postal fiscal y su Certificado de Sello Digital (archivos \texttt{.cer} y \texttt{.key} con contraseña), y su contador confirma por escrito el tratamiento de IVA e IEPS y las claves de producto del SAT. Si faltan, la facturación se entrega probada en modo de pruebas y se activa cuando lleguen.
    \item \textbf{Equipo para el POS:} el cliente aporta una tableta o teléfono con navegador actualizado y conexión a internet.
    \item \textbf{Retroalimentación:} el cliente revisa cada entrega quincenal por WhatsApp en un plazo máximo de 3 días hábiles. Sin observaciones en ese plazo, la entrega se considera aceptada.
\end{itemize}

% ============================================================
% SECCIÓN 2: DESCRIPCIÓN GENERAL DEL SISTEMA
% ============================================================
\section{Descripción General del Sistema}

\subsection{Perspectiva del Producto y Referencias Visuales}
El sistema es una aplicación web con \textbf{9 páginas públicas} (\texttt{index}, \texttt{historia}, \texttt{proceso}, \texttt{mezcales}, \texttt{contacto} y una página por cada una de las 4 expresiones) y \textbf{2 páginas privadas} (\texttt{login} y \texttt{dashboard}), conectadas a un backend propio. Su diseño se inspira en la interactividad y narrativa visual de \textbf{400 Conejos} (\texttt{400conejos.mx}) y \textbf{Casa Chagoya} (\texttt{casachagoya.com}).
\begin{itemize}[leftmargin=1.5em]
    \item \textbf{Personalidad de marca:} Artesanal y auténtica, tradicional y familiar, sostenible y cercana a la tierra.
    \item \textbf{Paleta cromática:} Tonos naturales (tierra, verde agave, madera, carbón orgánico y acentos dorados/ámbar), elegida por el cliente en la pregunta 4.5.
    \item \textbf{Tono de comunicación propuesto:} Formal, sofisticado y con trasfondo de herencia familiar. Pendiente de confirmar en la pregunta 4.6.
    \item \textbf{Restricciones de diseño explícitas del cliente:}
    \begin{itemize}
        \item Prohibido el aspecto de ``presentación de PowerPoint'' o diseño plano/simple.
        \item Prohibidas las animaciones rotas o bruscas (``que las animaciones no agarren'').
        \item Prohibidas las imágenes sobrepuestas que no se vean integradas orgánicamente con el fondo del sitio.
    \end{itemize}
\end{itemize}

\subsection{Usuarios, Roles y Permisos}
El sistema distingue cuatro tipos de usuario. Los permisos se validan en el servidor, no solo en la interfaz.

\renewcommand{\arraystretch}{1.3}
\begin{longtable}{@{} >{\raggedright\arraybackslash}p{6.0cm} >{\centering\arraybackslash}p{1.7cm} >{\centering\arraybackslash}p{1.6cm} >{\centering\arraybackslash}p{1.9cm} >{\centering\arraybackslash}p{2.6cm} @{}}
\toprule
\textbf{Función} & \textbf{Visitante} & \textbf{Cliente} & \textbf{Vendedor POS} & \textbf{Administrador} \\
\midrule
\endhead
\bottomrule
\endfoot
Navegar el sitio, fichas, galería y descargar PDF & \cmark & \cmark & \cmark & \cmark \\
Usar el carrito, enviar un pedido y pagar en línea & \cmark & \cmark & --- & --- \\
Solicitar y descargar la factura de un pedido propio & \cmark & \cmark & --- & \cmark \\
Consultar seguimiento con el enlace del correo de confirmación & \cmark & \cmark & --- & \cmark \\
Ver historial de pedidos propios en su cuenta & --- & \cmark & --- & --- \\
Usar el chat en vivo como visitante & \cmark & \cmark & --- & --- \\
Responder chats desde la bandeja del dashboard & --- & --- & \cmark & \cmark \\
Registrar ventas presenciales en el POS & --- & --- & \cmark & \cmark \\
Confirmar pagos y cambiar el estado de los pedidos & --- & --- & --- & \cmark \\
Gestionar productos, precios, lotes e inventario & --- & --- & --- & \cmark \\
Ver métricas de ventas y exportar reportes & --- & --- & --- & \cmark \\
Emitir reembolsos, cancelar facturas y emitir la factura global & --- & --- & --- & \cmark \\
Crear usuarios Vendedor POS y Administrador & --- & --- & --- & \cmark \\
\end{longtable}

\subsection{Arquitectura y Stack Tecnológico}
\renewcommand{\arraystretch}{1.3}
\begin{longtable}{@{} >{\raggedright\arraybackslash}p{3.6cm} >{\raggedright\arraybackslash}p{12.2cm} @{}}
\toprule
\textbf{Capa} & \textbf{Tecnología y función} \\
\midrule
\endhead
\bottomrule
\endfoot
Frontend & HTML5, Tailwind CSS compilado para producción y JavaScript nativo, a partir del prototipo actual. \\
Backend & API REST en \textbf{Go} (última versión estable), compilada como un solo binario. Incluye los webhooks de WhatsApp y de la pasarela de pago, y el servidor WebSocket del chat. \\
Base de datos & PostgreSQL en el mismo servidor, con migraciones versionadas. \\
Servidor web y HTTPS & Caddy como proxy inverso con certificados Let's Encrypt renovados automáticamente. \\
Hospedaje & VPS Vultr en la Ciudad de México: 1 vCPU, 1~GB de RAM y 25~GB SSD. Los datos permanecen en México. \\
DNS, CDN y protección & Cloudflare, plan gratuito. \\
Almacenamiento de archivos & Cloudflare R2 para videos, fichas PDF y respaldos diarios de la base de datos. \\
Correo electrónico & Resend para el envío de correos y Cloudflare Email Routing para recibir \textit{contacto@} y \textit{ventas@}. \\
WhatsApp & WhatsApp Business Platform (Cloud API oficial de Meta). \\
Pagos en línea & Mercado Pago Checkout Pro (o Stripe como alternativa): tarjeta, OXXO y SPEI. Los datos de tarjeta no pasan por el servidor del sitio. \\
Facturación & API de un PAC autorizado por el SAT (Facturama o Facturapi) para timbrar CFDI 4.0. \\
\end{longtable}

\subsection{Restricciones Operativas y de Capacidad}
Estas restricciones se derivan del presupuesto de infraestructura de \$1,550.00 MXN, de los planes gratuitos elegidos y de los servicios de pago por uso:
\begin{itemize}[leftmargin=1.5em]
    \item \textbf{Correos:} el plan gratuito de Resend permite 3,000 correos al mes y 100 al día, sumando transaccionales y newsletter. Un boletín para más de 100 suscriptores se envía en varios días o requiere el plan Pro de US\$20.00 al mes.
    \item \textbf{Archivos:} el almacenamiento gratuito de R2 es de 10~GB. Los videos se publican optimizados a 720p para respetar ese límite.
    \item \textbf{Servidor:} el plan de 1~GB está dimensionado para el volumen inicial. Si el uso de CPU o RAM supera el 80\% de forma sostenida, se recomienda escalar al plan de 2~GB.
    \item \textbf{WhatsApp:} los primeros 1,000 mensajes de servicio de cada mes son gratuitos. Los mensajes adicionales y las plantillas de marketing tienen costo por mensaje, a cargo del cliente.
    \item \textbf{Pagos en línea:} cada cobro paga una comisión de la pasarela, de alrededor de 3\% a 4\% más una cuota fija y el IVA. Los pagos en OXXO se acreditan hasta que el cliente paga en tienda y vencen a los 3 días.
    \item \textbf{Facturación:} cada factura timbrada cuesta entre \$0.50 y \$2.00 MXN según el PAC. Las facturas se solicitan dentro del mes de la compra; después se incluyen en la factura global.
\end{itemize}

% ============================================================
% SECCIÓN 3: REQUERIMIENTOS FUNCIONALES (RF)
% ============================================================
\section{Requerimientos Funcionales (RF)}

Todos los requerimientos de esta sección son \textbf{Indispensables para la Fase 1}. La columna \textit{Estado en Prototipo} refleja la revisión del código del prototipo realizada el 26 de octubre de 2026. ``Simulado'' significa que la interfaz existe pero no guarda ni envía datos reales.

\newcolumntype{L}[1]{>{\raggedright\arraybackslash}p{#1}}
\renewcommand{\arraystretch}{1.35}
\small
\begin{longtable}{@{} L{1.2cm} L{3.0cm} L{7.4cm} L{3.9cm} @{}}
\toprule
\textbf{ID} & \textbf{Módulo / Sección} & \textbf{Descripción del Requerimiento} & \textbf{Estado en Prototipo} \\
\midrule
\endfirsthead
\toprule
\textbf{ID} & \textbf{Módulo / Sección} & \textbf{Descripción del Requerimiento} & \textbf{Estado en Prototipo} \\
\midrule
\endhead
\bottomrule
\endfoot

\multicolumn{4}{@{}l}{\textbf{\textcolor{mezcalEarth}{Bloque A. Sitio institucional y de marca}}} \\
\midrule
\textbf{RF-01} & Verificación de Edad (+18) & Al ingresar al sitio, bloquear el scroll y difuminar el contenido mostrando un modal de mayoría de edad. Guardar la confirmación en el navegador. Mostrar un mensaje de restricción si el usuario selecciona ``No''. & Implementado \\
\midrule
\textbf{RF-02} & Internacionalización (ES / EN) & Selector de idioma en la barra de navegación (escritorio y móvil) que traduzca el \textbf{100\% de los textos públicos}: fichas, proceso, formularios, botones, correos al cliente y avisos legales. El dashboard y el POS son de uso interno y solo en español. & Parcial: 5 páginas traducidas; faltan las 4 páginas individuales de producto \\
\midrule
\textbf{RF-03} & Hero / Inicio Impactante & Encabezado a pantalla completa con video de fondo o video \textit{loopable} de la botella generado con IA, animación tipográfica de \textit{``MIL SENSACIONES''}, tagline e indicador de scroll. & Parcial: video del cliente integrado; faltan los loops de botella y el tagline oficial \\
\midrule
\textbf{RF-04} & Historia y Valores de la Marca & Sección narrativa sobre las 3 generaciones (abuelo, padre e hijo), origen en \textit{El Gramal, Nejapa de Madero, Yautepec, Oaxaca}, misión y visión de la marca. & Parcial: falta Misión y Visión del cliente \\
\midrule
\textbf{RF-05} & Perfil del Maestro Mezcalero & Sección dedicada a \textbf{Damián Olivera Díaz} (2.ª generación de maestros mezcaleros), su trayectoria y fotografía real integrada al estilo visual. & Parcial: usa una foto de Instagram; falta foto en alta resolución \\
\midrule
\textbf{RF-06} & Fichas Detalladas de Productos & Página propia para cada expresión: \textbf{Espadín 750~ml} (actual), \textbf{Ensamble}, \textbf{Silvestre} y \textbf{Curados}. Cada ficha muestra los \textbf{10 atributos} de la Sección~3.2 y, si el cliente lo decide, la etiqueta ``Próximamente''. & Parcial: 4 páginas creadas con datos provisionales por confirmar \\
\midrule
\textbf{RF-07} & Ficha Técnica en PDF & Botón en cada producto para ver y descargar su ficha técnica oficial en PDF, en español e inglés. & Parcial: ficha imprimible desde el navegador; falta el PDF descargable \\
\midrule
\textbf{RF-08} & Proceso de Elaboración & Explicación interactiva del proceso artesanal (selección y corte, cocción en horno cónico, molienda, fermentación, destilación y envasado) con fotos o video real. & Implementado; falta validar los textos con el cliente \\
\midrule
\textbf{RF-09} & Galería Multimedia & Galería interactiva con \textit{lightbox} de fotografías y video del palenque, campos de maguey y destilación. & Implementado \\
\midrule
\textbf{RF-10} & Dónde Comprar / Google Maps & Mapa embebido de Google Maps con el pin exacto del palenque, más la lista de puntos de venta autorizados. & Parcial: mapa embebido; faltan pin exacto y puntos de venta \\
\midrule
\textbf{RF-11} & Formulario de Contacto & Formulario con validación (Nombre, Correo, Teléfono/Empresa, Tipo de consulta: Consumidor / Bar-Restaurante / Importador, Mensaje). Guarda la consulta y la envía por correo a \textit{contacto@}. & Simulado \\
\midrule
\textbf{RF-12} & Enlaces a WhatsApp & Botón flotante y enlaces \texttt{wa.me} con el número real de la marca y mensaje prellenado según el producto o la sección. & Parcial: número ficticio \texttt{5219510000000} \\
\midrule
\textbf{RF-13} & Suscripción a Newsletter & Captura de correo con doble confirmación para el boletín de lotes y lanzamientos. Cada correo incluye un enlace para darse de baja. & Simulado \\
\midrule
\textbf{RF-14} & Redes Sociales y Compartir & Enlaces a Instagram (\texttt{@mezcalmilsensaciones}) y Facebook oficial, más botones para compartir fichas o el sitio. & Implementado; faltan botones en las páginas individuales \\
\midrule
\textbf{RF-15} & Aviso de Privacidad y Legal & Aviso de Privacidad conforme a la LFPDPPP, Términos y Condiciones de Venta conforme a la LFPC (envíos, devoluciones por daño, venta solo a mayores de 18), aviso de cookies y leyenda sanitaria obligatoria en todas las páginas públicas. & Parcial: aviso de privacidad solo en inicio; faltan Términos y leyenda en 4 páginas \\
\midrule
\multicolumn{4}{@{}l}{\textbf{\textcolor{mezcalEarth}{Bloque B. Plataforma comercial y de gestión}}} \\
\midrule
\textbf{RF-16} & Carrito de Compras Slidebar & Panel lateral para agregar productos, cambiar cantidades, eliminar y ver subtotal y total en MXN. Conserva el carrito al recargar y lo asocia a la cuenta al iniciar sesión. & Parcial: funciona en el navegador; las páginas de producto cargan el carrito desde una ruta local que fallará en producción \\
\midrule
\textbf{RF-17} & Registro de Pedido & El cliente captura nombre, teléfono, correo y dirección de envío. El sistema guarda el pedido con folio único y estado \textit{Orden Recibida}, envía la confirmación por correo y lo dirige al pago en línea (RF-26). Si elige pago manual, abre WhatsApp con el resumen y el folio. & No implementado \\
\midrule
\textbf{RF-18} & Cuentas de Usuario & Registro con correo y contraseña, verificación del correo, confirmación de mayoría de edad, inicio y cierre de sesión, y recuperación de contraseña con enlace temporal. & Simulado: credenciales de prueba fijas en \texttt{login.html} \\
\midrule
\textbf{RF-19} & Control de Acceso por Rol & Aplicar la matriz de la Sección~2.2. El rol lo asigna el Administrador, nunca el usuario al iniciar sesión. & Simulado: el rol se elige en un selector \\
\midrule
\textbf{RF-20} & Dashboard Administrativo & Alta y edición de productos y precios. Registro de lotes (número, agave, fecha de destilación, kilos de maguey en horno, \% Alc. Vol., botellas producidas, estado de marbetes CRM). Inventario por lote, gestión de pedidos y clientes, métricas de ingresos por mes y por producto, y exportación a CSV. & Simulado: datos de ejemplo \\
\midrule
\textbf{RF-21} & Terminal Punto de Venta (POS) & Pantalla táctil para ventas presenciales: búsqueda de producto, cantidad, método de pago (efectivo, transferencia o terminal externa), registro de la venta, descuento de inventario del lote y ticket digital por correo o WhatsApp. & Simulado \\
\midrule
\textbf{RF-22} & Seguimiento de Pedidos & Estados: \textit{Orden Recibida, Pago Pendiente, Pago Confirmado, En Empaque Artesanal, En Tránsito, Entregado} o \textit{Cancelado}. El Administrador captura paquetería y número de guía. El cliente consulta el estado en su cuenta o con el enlace del correo y recibe un correo en cada cambio. Cada pedido queda vinculado al lote de origen. & No implementado \\
\midrule
\textbf{RF-23} & Chatbot de WhatsApp & Menú automático 24/7 con opciones: catálogo y precios, mayoreo para bares y restaurantes, atención a importadores, estado de pedido por folio y hablar con una persona. Fuera del horario de atención, responde con un mensaje automático. & No implementado \\
\midrule
\textbf{RF-24} & Chat en Vivo en el Sitio & Ventana flotante en las páginas públicas. El visitante deja nombre y correo y conversa en tiempo real. Vendedores y Administradores responden desde una bandeja en el dashboard con historial. Si nadie está conectado, el mensaje se guarda y se notifica por correo. & No implementado \\
\midrule
\textbf{RF-25} & Correos Transaccionales & Envío automático de bienvenida, verificación de cuenta, recuperación de contraseña, confirmación de pedido y de pago, cambio de estado, ticket del POS, envío de factura y boletín. El dominio queda autenticado con SPF, DKIM y DMARC. & No implementado \\
\midrule
\textbf{RF-26} & Pagos en Línea & Pago con tarjeta de crédito o débito, efectivo en OXXO o transferencia SPEI mediante el checkout de la pasarela. El webhook confirma el pago y cambia el pedido a \textit{Pago Confirmado}. Un pago en OXXO deja el pedido en \textit{Pago Pendiente} y aparta el inventario hasta su vencimiento. El Administrador emite reembolsos totales o parciales desde el dashboard. & No implementado \\
\midrule
\textbf{RF-27} & Facturación Electrónica CFDI 4.0 & El cliente solicita la factura desde su cuenta o con el folio, captura RFC, razón social, régimen fiscal, código postal y uso del CFDI, y recibe PDF y XML por correo. El sistema timbra con el PAC usando el CSD del negocio y desglosa IVA e IEPS según la configuración del contador. El Administrador consulta, reenvía y cancela facturas y emite la factura global a público en general. & No implementado \\
\end{longtable}
\normalsize

\subsection{Criterios de Aceptación}
Cada requerimiento se da por cumplido cuando su criterio se verifica en el entorno de producción.

\renewcommand{\arraystretch}{1.3}
\small
\begin{longtable}{@{} L{1.9cm} L{13.9cm} @{}}
\toprule
\textbf{ID} & \textbf{Criterio de aceptación verificable} \\
\midrule
\endfirsthead
\toprule
\textbf{ID} & \textbf{Criterio de aceptación verificable} \\
\midrule
\endhead
\bottomrule
\endfoot
RF-01 & Al abrir cualquier página pública por primera vez aparece el modal. Con ``No'' el contenido sigue bloqueado. Con ``Sí'' no vuelve a aparecer en ese navegador. \\
RF-02 & Al cambiar a inglés, ningún texto visible de las 9 páginas públicas queda en español, y el idioma elegido se conserva al navegar. \\
RF-03 & El hero carga en móvil y escritorio sin saltos visibles, y el video reproduce en bucle continuo sin audio. \\
RF-04 / RF-05 & Las secciones muestran los textos y fotografías oficiales entregados por el cliente, sin imágenes de relleno. \\
RF-06 & Las 4 fichas muestran los 10 atributos con datos confirmados por escrito por el cliente. \\
RF-07 & Cada ficha descarga un PDF en el idioma activo con los 10 atributos y el logo. \\
RF-08 / RF-09 & Las 6 etapas del proceso y la galería abren fotos y videos en \textit{lightbox} y se cierran con la tecla Esc o un toque fuera. \\
RF-10 & El mapa abre el pin exacto del palenque y la lista de puntos de venta coincide con la entregada por el cliente. \\
RF-11 & Un envío de prueba llega a \textit{contacto@} en menos de 2 minutos y queda guardado en la base de datos. Los campos inválidos muestran un mensaje de error. \\
RF-12 & Todos los enlaces abren WhatsApp con el número oficial y un mensaje prellenado acorde a la página. \\
RF-13 & Un correo de prueba recibe la confirmación. Tras confirmar queda suscrito, y el enlace de baja lo elimina de la lista. \\
RF-14 & Los íconos abren los perfiles oficiales y el botón compartir funciona en Android, iOS y escritorio. \\
RF-15 & Las 9 páginas públicas enlazan el Aviso de Privacidad y los Términos, y muestran la leyenda sanitaria. \\
RF-16 & El carrito conserva productos y cantidades al recargar, y el total coincide con la suma de precio por cantidad. \\
RF-17 & Un pedido de prueba genera folio, aparece en el dashboard, envía el correo de confirmación y lleva al pago en línea. Con pago manual, abre WhatsApp con el resumen. \\
RF-18 & Un usuario nuevo se registra, verifica su correo, inicia sesión y recupera su contraseña. Un enlace de recuperación usado o vencido es rechazado. \\
RF-19 & Un Cliente o Vendedor POS que intenta abrir una función de Administrador recibe un error de acceso, aun escribiendo la dirección directamente. \\
RF-20 & Un producto o lote creado en el dashboard aparece en el sitio sin modificar código, y el inventario coincide con las ventas registradas. \\
RF-21 & Una venta en el POS descuenta el inventario del lote, emite el ticket digital y aparece en las métricas del día. \\
RF-22 & Cada cambio de estado envía un correo al cliente y se refleja en su cuenta y en el enlace de seguimiento. \\
RF-23 & Cada opción del menú responde correctamente desde un teléfono externo, y la consulta por folio devuelve el estado real del pedido. \\
RF-24 & Un mensaje enviado desde el sitio aparece en la bandeja del dashboard en menos de 3 segundos, y la respuesta llega al visitante. \\
RF-25 & Los correos de prueba llegan a la bandeja principal de Gmail y Outlook, no a spam, y pasan la validación SPF, DKIM y DMARC. \\
RF-26 & Un pago real con tarjeta y uno en OXXO, ambos de bajo monto, cambian el pedido al estado correcto sin intervención manual. Un pedido de OXXO vencido libera el inventario. Un reembolso desde el dashboard se refleja en la pasarela. \\
RF-27 & Una factura de prueba se timbra, llega por correo con PDF y XML, y aparece como vigente en el verificador de CFDI del SAT. Una cancelación desde el dashboard se refleja en el SAT. \\
\end{longtable}
\normalsize

\subsection{Detalle Obligatorio de la Ficha de Producto (10 Atributos)}
Según la Sección 2.4 del cuestionario, cada producto presentado en el sitio debe mostrar obligatoriamente estos \textbf{10 campos}:
\begin{enumerate}[leftmargin=1.8em]
    \item \textbf{Nombre del producto} (ej. \textit{Mil Sensaciones — Espadín Joven}).
    \item \textbf{Tipo de agave utilizado} (nombre común y científico, ej. \textit{Agave Angustifolia Haw}, tiempo de maduración).
    \item \textbf{Maestro Mezcalero} (\textit{Damián Olivera Díaz}).
    \item \textbf{Región y comunidad de origen} (\textit{El Gramal, Nejapa de Madero, Yautepec, Oaxaca}).
    \item \textbf{Proceso de elaboración} (cocción, molienda, fermentación, tipo de alambique y número de destilaciones).
    \item \textbf{Graduación alcohólica} (\texttt{\% Alc. Vol.} real por lote).
    \item \textbf{Notas de cata} (vista, aromas en nariz y sabores en boca).
    \item \textbf{Presentaciones disponibles} (\texttt{750 ml} u otras capacidades).
    \item \textbf{Cantidad de botellas por lote} (ej. \textit{Edición limitada de $X$ botellas por lote}).
    \item \textbf{Código NOM / Certificación CRM} (NOM-070-SCFI-2016).
\end{enumerate}

\textbf{Importante:} los valores que hoy muestra el prototipo (45\%, 47\%, 48\% y 38\% Alc. Vol., micro-lote de 300 botellas, notas de cata y maridajes) son \textbf{provisionales}. Se reemplazan por los datos que el cliente entregue en la Sección~6, Parte~C.

% ============================================================
% SECCIÓN 4: REQUERIMIENTOS NO FUNCIONALES (RNF)
% ============================================================
\section{Requerimientos No Funcionales (RNF)}
\begin{itemize}[leftmargin=1.5em]
    \item \textbf{RNF-01 Rendimiento y Fluidez Visual:} en móvil con red 4G, LCP menor a 2.5~s, INP menor a 200~ms y CLS menor a 0.1, medidos con PageSpeed Insights. Las animaciones corren a 60~FPS usando \texttt{transform} y \texttt{opacity}. Las respuestas de la API para catálogo y pedidos tardan menos de 300~ms en el 95\% de los casos.
    \item \textbf{RNF-02 Diseño Responsivo (Mobile-First):} adaptación completa a teléfonos (320--767~px), tabletas (768--1023~px) y escritorio (1024~px o más), incluyendo selector de idioma, carrito, chat y POS.
    \item \textbf{RNF-03 Optimización Multimedia:} imágenes en \texttt{WebP} o \texttt{AVIF} con carga diferida. Las botellas se integran como imagen sin fondo o como video \textit{loopable} 720p sin audio sobre un fondo del mismo color del sitio. Cada video de loop pesa como máximo 5~MB.
    \item \textbf{RNF-04 SEO y Metadatos:} etiquetas \texttt{<meta>} y Open Graph bilingües, datos estructurados de Schema.org para producto y organización, \texttt{sitemap.xml}, \texttt{robots.txt} y etiquetas \texttt{hreflang} para español e inglés.
    \item \textbf{RNF-05 Mantenibilidad:} productos, precios y lotes se actualizan desde el dashboard sin tocar código. Los textos institucionales los actualiza el equipo de desarrollo, sin CMS en la Fase 1. El código incluye migraciones de base de datos versionadas, pruebas automáticas de los flujos de registro, pedido, pago, facturación e inventario, y un manual de despliegue.
    \item \textbf{RNF-06 Seguridad:} HTTPS obligatorio con HSTS. Contraseñas almacenadas con Argon2id. Sesión con JWT en cookie \texttt{HttpOnly}, \texttt{Secure} y \texttt{SameSite}. Bloqueo temporal tras 5 intentos fallidos de inicio de sesión en 15 minutos. Validación de datos en el servidor, consultas parametrizadas y protección contra las vulnerabilidades del OWASP Top 10. Los webhooks de la pasarela se aceptan solo con firma válida. El sitio no guarda datos de tarjetas, y el CSD y su contraseña se almacenan cifrados. Las credenciales de prueba del prototipo se eliminan antes del lanzamiento.
    \item \textbf{RNF-07 Disponibilidad y Respaldos:} disponibilidad objetivo de 99.5\% mensual, sin contar mantenimientos avisados con anticipación. Respaldo automático diario de la base de datos en Cloudflare R2, con retención de 30 días. Pérdida máxima aceptable de 24 horas de datos y restauración en menos de 4 horas. Se realiza una prueba de restauración antes del lanzamiento.
    \item \textbf{RNF-08 Privacidad de Datos:} se recaban solo los datos necesarios para vender y entregar. Los datos se alojan en México. El Aviso de Privacidad indica el correo para ejercer los derechos ARCO. Sin consentimiento solo se usan cookies técnicas.
    \item \textbf{RNF-09 Compatibilidad:} últimas dos versiones de Chrome, Safari en iOS, Firefox, Edge y Samsung Internet.
    \item \textbf{RNF-10 Accesibilidad:} contraste de texto nivel AA de WCAG 2.2, navegación con teclado en menús y modales, texto alternativo en imágenes y animaciones reducidas cuando el usuario activa la preferencia del sistema \texttt{prefers-reduced-motion}.
    \item \textbf{RNF-11 Cumplimiento para Bebidas Alcohólicas:} verificación de mayoría de edad al entrar, al crear cuenta y al enviar pedido. Leyenda sanitaria visible en todas las páginas públicas. Código NOM y certificación CRM visibles en cada ficha. La publicidad y el boletín se dirigen solo a mayores de 18 años.
\end{itemize}

\newpage

% ============================================================
% SECCIÓN 5: REGISTRO ÍNTEGRO DEL CUESTIONARIO DE LEVANTAMIENTO
% ============================================================
\section{Registro del Cuestionario de Levantamiento (Cuestionario Original)}

A continuación se documentan las respuestas entregadas por el cliente en el cuestionario original, destacando también aquellas preguntas que quedaron \textbf{sin responder o incompletas} (las cuales tienen su espacio de respuesta en la Sección 6).

\subsection{Sección 1 — La Marca}
\begin{itemize}[leftmargin=1.5em]
    \item \textbf{1.1 ¿Cuál es el nombre de la marca de mezcal?}\\
    \textbf{Respuesta:} \texttt{MIL SENSACIONES}

    \item \textbf{1.2 ¿Dónde producen el mezcal? (estado, municipio, comunidad)}\\
    \textbf{Respuesta:} El Gramal, Nejapa de Madero, Yautepec, Oaxaca.

    \item \textbf{1.3 ¿Cuánto tiempo llevan produciendo mezcal?}\\
    \textbf{Respuesta:} Se empezó a producir nuevamente hace dos años.

    \item \textbf{1.4 ¿Qué hace diferente a su mezcal del resto del mercado?}
    \begin{answerbox}
    ``Nuestro mezcal Mil Sensaciones no es solo una bebida\dots es historia viva. Somos la tercera generación de una tradición con el abuelo, continuó con el padre y hoy sigue en manos del hijo. Cada gota lleva años de experiencia, dedicación y respeto por el proceso artesanal.\\
    Aquí no improvisamos. Conservamos las enseñanzas originales, los tiempos, los métodos y el cuidado que le dan a nuestro mezcal ese sabor único, auténtico e inconfundible.\\
    Seleccionamos cuidadosamente los mejores magueyes, respetamos sus ciclos y trabajamos cada etapa con pasión: desde el corte hasta la destilación.\\
    No buscamos producir masa, buscamos crear algo que realmente se sienta.\\
    Mil Sensaciones tiene carácter, tiene identidad\dots y sobre todo, tiene alma.\\
    Porque no vendemos solo mezcal, compartimos una herencia, una cultura y una experiencia que se disfruta en cada trago.''
    \end{answerbox}

    \item \textbf{1.5 ¿Tienen un maestro mezcalero identificado? ¿Quién es y cuál es su historia?}
    \begin{answerbox}
    ``Damián Olivera Díaz es el corazón de Mezcal Mil Sensaciones.\\
    Segunda generación de maestros mezcaleros, su historia comienza desde la niñez, entre magueyes, aprendiendo a sembrar, cuidar y respetar la tierra. Lo que para muchos es un proceso, para él ha sido una forma de vida.\\
    Con los años, perfeccionó su técnica sin perder la esencia: hacer mezcal como se hacía antes, con paciencia, conocimiento y pasión.\\
    Hoy juntos, no solo elaboramos mezcal\dots creamos emociones en cada botella.\\
    Porque detrás de cada gota, hay historia, familia y alma.''
    \end{answerbox}

    \item \textbf{1.6 ¿Cuál es la historia detrás de la marca? (origen, familia, motivación para crearla)}
    \begin{answerbox}
    ``Mezcal Mil Sensaciones nace del legado de una familia que ha vivido el mezcal desde sus raíces. No comenzó como un negocio, sino como una tradición: el conocimiento transmitido de generación en generación, entre magueyes, tierra y fuego.\\
    Desde el abuelo hasta hoy, la motivación siempre ha sido la misma: preservar la esencia del mezcal auténtico y compartirlo con quienes saben valorar su origen.\\
    La marca surge del deseo de llevar esa herencia más allá, de transformar años de trabajo, aprendizaje y pasión en una experiencia que conecte con los sentidos.\\
    Mil Sensaciones no solo representa un mezcal\dots Representa historia, identidad y el orgullo de una familia que decidió compartir su esencia con el mundo.''
    \end{answerbox}

    \item \textbf{1.7 Elementos de la marca con los que ya cuentan (todos marcados por el cliente):}\\
    \chk\ Nombre de la marca definido \quad
    \chk\ Logo diseñado \quad
    \chk\ Colores de marca seleccionados \quad
    \chk\ Slogan o tagline \\
    \chk\ Misión y visión redactadas \quad
    \chk\ Etiqueta de botella diseñada \quad
    \chk\ Registro de marca (IMPI) \\
    \chk\ Certificación NOM-070 / CRM \quad
    \chk\ Redes sociales activas \quad
    \chk\ Correo electrónico de la empresa

    \item \textbf{1.8 Redes sociales activas:}
    \begin{itemize}
        \item \textbf{Instagram:} \url{https://www.instagram.com/mezcalmilsensaciones/}
        \item \textbf{Facebook:} \url{https://www.facebook.com/profile.php?id=61581942955039}
    \end{itemize}
\end{itemize}

\subsection{Sección 2 — Los Productos}
\begin{itemize}[leftmargin=1.5em]
    \item \textbf{2.1 ¿Cuántos productos tienen actualmente o planean tener al lanzar el sitio?}\\
    \textbf{Respuesta:} \textbf{750 ml — Mezcal Espadín} (Próximamente a lanzar: \textit{Ensambles}, \textit{Mezcal Silvestre}, \textit{Curados de mezcal}).

    \item \textbf{2.2 Tipos de mezcal que producen o producirán:}\\
    \chk\ Mezcal Joven \quad
    \uchk\ Mezcal Reposado \quad
    \uchk\ Mezcal Añejo \quad
    \chk\ Edición Limitada \\
    \chk\ Ensamble (mezcla de agaves) \quad
    \uchk\ Crema de Mezcal \quad
    \chk\ Otro: \textit{Curados de mezcal}.

    \item \textbf{2.4 Información que desean mostrar en la ficha de cada producto:}\\
    \chk\ Nombre del producto \quad
    \chk\ Tipo de agave utilizado \quad
    \chk\ Maestro mezcalero \quad
    \chk\ Región y comunidad de origen \\
    \chk\ Proceso de elaboración \quad
    \chk\ Graduación alcohólica (\% Alc. Vol.) \quad
    \chk\ Notas de cata (sabores y aromas) \\
    \uchk\ Maridajes sugeridos \quad
    \chk\ Presentaciones disponibles (ml) \quad
    \uchk\ Premios o reconocimientos \\
    \chk\ Cantidad de botellas por lote \quad
    \chk\ Código NOM \quad
    \uchk\ Precio de referencia

    \item \textbf{2.5 Material visual actual para los productos:}\\
    \chk\ Fotos básicas (teléfono), pero sirven de referencia.\\
    \chk\ Tenemos video del proceso o la palenque.\\
    \chk\ Fotos tomadas con cámara profesional por nosotros mismos.

    \item \textbf{2.6 ¿Tienen ya redactada la descripción o historia de alguno de sus productos?}\\
    \textcolor{red!80!black}{\textbf{[SIN RESPUESTA EN EL CUESTIONARIO — Ver espacio de llenado en Sección 6.1]}}
\end{itemize}

\subsection{Sección 3 — El Sitio Web}
\begin{itemize}[leftmargin=1.5em]
    \item \textbf{3.1 Secciones indispensables marcadas para la primera versión:}\\
    \chk\ Inicio (página principal con video o foto impactante) \quad
    \chk\ Historia de la marca / Nosotros \\
    \uchk\ Catálogo de productos (sin precio ni carrito) \quad
    \chk\ Ficha detallada de cada producto \\
    \chk\ Proceso de elaboración \quad
    \chk\ Maestro mezcalero (perfil y trayectoria) \\
    \uchk\ Dónde encontrarnos / puntos de venta \textit{(Nota: marcado como Necesario en 3.4 y 7.1)} \\
    \uchk\ Contacto y formulario de consulta \textit{(Nota: marcado como Necesario en 3.4 y 7.1)} \\
    \uchk\ Prensa / Kit de medios \quad
    \chk\ Galería de fotos \quad
    \uchk\ Blog o noticias \quad
    \uchk\ Aviso de privacidad

    \item \textbf{3.2 ¿El sitio necesita estar en inglés también?}\\
    \uradio\ No, solo en español por ahora\\
    \radio\ \underline{\textbf{Sí, español e inglés desde el inicio}} \quad \textcolor{mezcalAgave}{\textbf{[RESUELTO: incluido en la Fase 1, ver RF-02]}}\\
    \uradio\ Español primero, inglés en una segunda fase

    \item \textbf{3.3 ¿Necesitan un panel de administración para editar el contenido sin programador?}\\
    \uradio\ Sí, alguien del equipo lo actualizará constantemente\\
    \uradio\ No por ahora, el equipo de desarrollo hará los cambios\\
    \radio\ \underline{\textbf{No sé, depende del costo}} \quad \textcolor{mezcalAgave}{\textbf{[RESUELTO: Incluido en el paquete integral de \$7,550.00 MXN a 8 semanas con Dashboard Administrativo + POS, ver RF-20 y RF-21]}}\

    \item \textbf{3.4 Otras funciones a incluir en el lanzamiento:}\\
    \chk\ Verificación de edad (+18) al entrar al sitio \quad
    \chk\ Formulario de contacto con notificación por correo \\
    \chk\ Enlace a WhatsApp para consultas directas \quad
    \uchk\ Integración con Instagram (posts recientes) \\
    \chk\ Google Maps con ubicación de la palenque \quad
    \chk\ Suscripción a lista de correo / newsletter \\
    \chk\ Descarga de ficha técnica en PDF \quad
    \chk\ Botones para compartir en redes sociales

    \item \textbf{3.5 ¿Hay alguna función adicional no listada que sea importante?}\\
    \textcolor{red!80!black}{\textbf{[SIN RESPUESTA EN EL CUESTIONARIO — Ver espacio de llenado en Sección 6.1]}}
\end{itemize}

\subsection{Sección 4 — Diseño y Estilo Visual}
\begin{itemize}[leftmargin=1.5em]
    \item \textbf{4.1 Personalidad de la marca (hasta 3):}\\
    \chk\ Artesanal y auténtica \quad
    \chk\ Tradicional y familiar \quad
    \chk\ Sostenible y cercana a la tierra

    \item \textbf{4.2 ¿Qué les gusta específicamente de \texttt{400conejos.mx} y \texttt{casachagoya.com}?}\\
    \textbf{Respuesta:} ``Lo interactivo que están los sitios, la combinación de colores''.

    \item \textbf{4.3 ¿Otro sitio web de referencia visual?}\\
    \textbf{Respuesta:} No.

    \item \textbf{4.4 ¿Qué NO quieren que se vea en el diseño del sitio?}\\
    \textbf{Respuesta:} ``Que las animaciones no agarren, o que parezca presentación de PowerPoint, que esté simple, imágenes que no se vean integradas''.

    \item \textbf{4.5 ¿Con qué paleta de colores se identifica más su marca?}\\
    \uradio\ Tonos oscuros (negro, café, miel, ámbar) \qquad
    \radio\ \underline{\textbf{Tonos naturales (tierra, verde, madera)}} \quad \textcolor{mezcalAgave}{\textbf{[IMPLEMENTADO]}}\\
    \uradio\ Tonos claros y minimalistas (blanco, beige, crema) \qquad
    \uradio\ Tonos vibrantes (terracota, rojo, naranja)\\
    \uradio\ No lo tenemos definido aún

    \item \textbf{4.6 ¿Cuál debe ser el tono de los textos del sitio?}\\
    \textcolor{red!80!black}{\textbf{[SIN SELECCIÓN EN EL CUESTIONARIO — Ver espacio de llenado en Sección 6.1]}}
\end{itemize}

\subsection{Sección 5 — Mercado y Clientes}
\begin{itemize}[leftmargin=1.5em]
    \item \textbf{5.1 Público objetivo del sitio web:}\\
    \chk\ Consumidor final \quad
    \chk\ Bares, restaurantes y coctelería \quad
    \chk\ Importadores internacionales \quad
    \chk\ Turistas y visitantes de Oaxaca / México \\
    \uchk\ Distribuidores nacionales \quad
    \uchk\ Periodistas, bloggers y sommeliers \quad
    \uchk\ Tiendas especializadas

    \item \textbf{5.2 Mercado principal:} Oaxaca.
    \item \textbf{5.3 Distribución fuera de México:} Sí, Estados Unidos.
    \item \textbf{5.4 Descripción del cliente ideal:}\\
    \textbf{Respuesta parcial:} ``Mayores de 18 años,'' \textcolor{red!80!black}{\textbf{[RESPUESTA INCOMPLETA — Ver Sección 6.1]}}
\end{itemize}

\subsection{Secciones 6, 7 y 8 — Técnico, Prioridades, Contenido y Tiempos}
\begin{itemize}[leftmargin=1.5em]
    \item \textbf{6.1 ¿Quién actualizará el contenido después de entregado?}\\
    \chk\ El mismo equipo de desarrollo (\textbf{Ing. en Computación Jose Ramon Aragon Toledo}).
    \item \textbf{6.2 Herramientas conectadas requeridas:} \chk\ WhatsApp Business.
    \item \textbf{7.1 Matriz de Prioridades de Lanzamiento:}
    \begin{itemize}
        \item \textbf{Sí / Necesario desde el inicio:} Historia y valores de la marca, Ficha detallada de cada mezcal, Proceso de elaboración explicado, Perfil del maestro mezcalero, Formulario o botón de contacto, Mapa o lista de dónde comprar, Enlace a redes sociales, Versión en inglés, Verificación de edad (+18).
        \item \textbf{Más adelante (Fase 2):} Catálogo visual de productos, Blog o sección de noticias, Panel de administración (CMS).
    \end{itemize}
    \item \textbf{7.2 Funciones deseadas a futuro (Fase 2):}\\
    \chk\ Tienda en línea \quad
    \chk\ Reservas para catas o tours a la palenque \quad
    \chk\ Portal para distribuidores con precios especiales \quad
    \chk\ App móvil \quad
    \chk\ Versión en inglés del sitio \quad
    \chk\ Blog sobre mezcal y cultura \quad
    \chk\ Integración con marketplaces (Amazon, MercadoLibre).
    \item \textbf{7.3 Indicadores de éxito esperados:} ``Visitas constantes al sitio, incremento de ventas, más seguidores e interacción en redes sociales''.
    \item \textbf{8.1 Contenido disponible actualmente para entregar:}\\
    \chk\ Nombre y logo de la marca \quad
    \chk\ Textos sobre la historia \quad
    \chk\ Misión y visión escritas \\
    \chk\ Fotografías de la palenque o del proceso \quad
    \chk\ Video del proceso de elaboración \quad
    \chk\ Fichas técnicas de los productos \\
    \uchk\ Fotografías profesionales de productos \quad
    \uchk\ Etiquetas o arte gráfico de botellas.
    \item \textbf{8.2 Redacción de textos:} ``Nosotros y el equipo de desarrollo''.
    \item \textbf{8.3 Toma de fotografías:} ``Lo hará alguien del equipo''.
    \item \textbf{8.4 Fecha límite o evento clave:} ``Por el momento no''.
    \item \textbf{8.5 Frecuencia y canal de retroalimentación:} ``En WhatsApp cuando se pueda''.
\end{itemize}

\newpage

% ============================================================
% SECCIÓN 6: ESPACIOS PARA PREGUNTAS Y DATOS PENDIENTES
% ============================================================
\section{Formulario de Cierre de Levantamiento (Preguntas y Datos Pendientes)}

Esta sección está diseñada para ser completada directamente con el cliente (\textbf{Mezcal Mil Sensaciones}) con el fin de resolver las preguntas que quedaron en blanco en el cuestionario original y recopilar los datos exactos necesarios para finalizar la Fase 1.

\subsection{Parte A: Preguntas del Cuestionario Original que Quedaron sin Responder}

\begin{pendingbox}{Pregunta 2.6 — Descripción o historia individual de cada producto}
\textbf{¿Tienen ya redactada la descripción o historia de alguno de sus productos? Si sí, favor de escribirla o adjuntarla; si no, confirmar si el equipo de desarrollo (\textit{Ing. en Computación Jose Ramon Aragon Toledo}) debe redactarla con base en las notas de cata:}
\fillline
\fillline
\fillline
\end{pendingbox}

\begin{pendingbox}{Pregunta 3.2 — Confirmación del Idioma Inglés (RESUELTO)}
\textbf{Respuesta subrayada por el cliente en el cuestionario original:}\\[0.4em]
\uradio\ No, solo en español por ahora\\[0.25em]
\radio\ \underline{\textbf{Sí, español e inglés desde el inicio}} \quad \textcolor{mezcalAgave}{\textbf{(Incluido en la Fase 1. En el prototipo ya funciona en 5 páginas; faltan las 4 páginas individuales de producto)}}\\[0.25em]
\uradio\ Español primero, inglés en una segunda fase
\end{pendingbox}

\begin{pendingbox}{Pregunta 3.3 — Confirmación sobre Panel de Administración (RESUELTO)}
\textbf{Respuesta subrayada por el cliente en el cuestionario original:}\\[0.4em]
\uradio\ Sí, alguien del equipo lo actualizará constantemente\\[0.25em]
\uradio\ No por ahora, el equipo de desarrollo hará los cambios\\[0.25em]
\radio\ \underline{\textbf{No sé, depende del costo}}\\[0.35em]
\textbf{Acuerdo Técnico-Económico:} El Panel de Administración (\texttt{dashboard.html}) con gestión de lotes/inventario, control de pedidos, punto de venta (POS) e historial de clientes \textbf{ya se incluye sin costo adicional} dentro del plan de desarrollo integral de \textbf{\$7,550.00 MXN a 8 semanas} (\$1,550.00 de infraestructura y \$6,000.00 de honorarios) elaborado por el \textbf{Ing. en Computación Jose Ramon Aragon Toledo}.
\end{pendingbox}

\begin{pendingbox}{Pregunta 3.5 — Funciones adicionales importantes no listadas}
\textbf{¿Existe alguna función adicional que no apareció en las listas anteriores y que sea importante para el lanzamiento?}
\fillline
\fillline
\end{pendingbox}

\begin{pendingbox}{Pregunta 4.5 — Selección Oficial de Paleta de Colores (RESUELTO)}
\textbf{Respuesta subrayada por el cliente en el cuestionario original:}\\[0.4em]
\uradio\ Tonos oscuros (negro, café, miel, ámbar)\\[0.2em]
\radio\ \underline{\textbf{Tonos naturales (tierra, verde, madera)}} \quad \textcolor{mezcalAgave}{\textbf{(Aplicado en el Sistema de Diseño Global)}}\\[0.2em]
\uradio\ Tonos claros y minimalistas (blanco, beige, crema)\\[0.2em]
\uradio\ Tonos vibrantes (terracota, rojo, naranja)\\[0.4em]
\textbf{Tokens cromáticos aplicados en el sitio web (\texttt{assets/css/styles.css}):}
\begin{itemize}[leftmargin=1.5em, topsep=2pt, itemsep=1pt]
    \item \textbf{Madera horneada / Encino (\texttt{mezcal-dark} / \texttt{mezcal-wood}):} \texttt{\#18110B} y \texttt{\#261A11} (fondos principales de secciones, barra de navegación y cajas translúcidas sobre videos).
    \item \textbf{Tierra de Oaxaca (\texttt{mezcal-charcoal} / \texttt{mezcal-earth}):} \texttt{\#2D1E14} y \texttt{\#8B5A2B} (tarjetas de productos, bordes cálidos, acentos de barro y franjas de sección).
    \item \textbf{Verde Maguey (\texttt{mezcal-agave} / \texttt{mezcal-agave-light}):} \texttt{\#4B6043} y \texttt{\#6C8663} (botones principales, insignias de origen, divisores botánicos y degradados orgánicos).
\end{itemize}
\end{pendingbox}

\begin{pendingbox}{Pregunta 4.6 — Tono de Comunicación de los Textos}
\textbf{Esta pregunta quedó sin marcar en el documento Word. Seleccione el tono deseado para la redacción del sitio:}\\[0.4em]
\uradio\ Formal y sofisticado \qquad
\uradio\ Cercano y conversacional\\[0.25em]
\uradio\ Poético y narrativo (que cuente historias) \qquad
\uradio\ Técnico y educativo (énfasis en el proceso)
\end{pendingbox}

\begin{pendingbox}{Pregunta 5.4 — Perfil Completo del Cliente Ideal}
\textbf{En el cuestionario únicamente se escribió \textit{``Mayores de 18 años,''}. Favor de completar: rango de edad principal, intereses, nivel socioeconómico y qué busca en un mezcal artesanal:}
\fillline
\fillline
\fillline
\end{pendingbox}

\newpage

\subsection{Parte B: Recolección de Textos Corporativos y Registros Oficiales}

En los puntos 1.7 y 8.1 del cuestionario el cliente indicó que \textbf{ya cuenta} con estos elementos, pero falta su texto/dato exacto para integrarlos en el código del sitio web:

\begin{pendingbox}{1. Identidad Verbal y Registros Regulatorios}
\textbf{a) Slogan o Tagline oficial de ``Mil Sensaciones'':}
\fillline
\fillline

\textbf{b) Misión de la empresa (texto íntegro):}
\fillline
\fillline
\fillline

\textbf{c) Visión de la empresa (texto íntegro):}
\fillline
\fillline
\fillline

\textbf{d) Código de Certificación NOM-070 / CRM exacto:} \rule{7.5cm}{0.5pt}\\[0.6em]
\textbf{e) Número de Registro de Marca (IMPI):} \rule{8.8cm}{0.5pt}
\end{pendingbox}

\subsection{Parte C: Fichas Técnicas Completas por Producto (10 Atributos Obligatorios)}

Complete los datos reales de cada expresión para generar sus tarjetas interactivas y sus respectivos archivos PDF descargables:

\renewcommand{\arraystretch}{1.65}
\begin{longtable}{@{} L{4.2cm} L{5.5cm} L{5.5cm} @{}}
\toprule
\textbf{Atributo Obligatorio} & \textbf{Producto 1: Mezcal Espadín (Actual)} & \textbf{Producto 2: Ensamble (Próximo)} \\
\midrule
\endfirsthead
\toprule
\textbf{Atributo Obligatorio} & \textbf{Producto 1: Mezcal Espadín (Actual)} & \textbf{Producto 2: Ensamble (Próximo)} \\
\midrule
\endhead
\bottomrule
\endfoot
1. Nombre comercial exacto & & \\ \hline
2. Tipo de agave (común, científico y años de maduración) & \textit{Agave Angustifolia Haw} (\rule{1.2cm}{0.4pt} años) & \\ \hline
3. Maestro Mezcalero & Damián Olivera Díaz & Damián Olivera Díaz \\ \hline
4. Región y comunidad & El Gramal, Nejapa de Madero, Yautepec, Oax. & El Gramal, Nejapa de Madero, Yautepec, Oax. \\ \hline
5. Proceso (horno, molienda, tina de fermentación, alambique) & & \\ \hline
6. Graduación alcohólica (\texttt{\% Alc. Vol.}) & \rule{2cm}{0.4pt} \% Alc. Vol. & \rule{2cm}{0.4pt} \% Alc. Vol. \\ \hline
7. Notas de cata (Vista, Nariz, Boca) & & \\ \hline
8. Presentaciones disponibles & 750 ml / Otra: \rule{2cm}{0.4pt} & 750 ml / Otra: \rule{2cm}{0.4pt} \\ \hline
9. Cantidad de botellas por lote & \rule{2.5cm}{0.4pt} botellas & \rule{2.5cm}{0.4pt} botellas \\ \hline
10. Código NOM / Lote & NOM: \rule{3cm}{0.4pt} & NOM: \rule{3cm}{0.4pt} \\
\end{longtable}

\begin{longtable}{@{} L{4.2cm} L{5.5cm} L{5.5cm} @{}}
\toprule
\textbf{Atributo Obligatorio} & \textbf{Producto 3: Mezcal Silvestre} & \textbf{Producto 4: Curados de Mezcal} \\
\midrule
1. Nombre comercial / Sabores & & Sabores: \\ \hline
2. Tipo de agave silvestre / Frutas o ingredientes & & \\ \hline
3. Maestro Mezcalero & Damián Olivera Díaz & Damián Olivera Díaz \\ \hline
4. Proceso de elaboración & & \\ \hline
5. Graduación alcohólica (\texttt{\% Alc. Vol.}) & \rule{2cm}{0.4pt} \% Alc. Vol. & \rule{2cm}{0.4pt} \% Alc. Vol. \\ \hline
6. Notas de cata (Vista, Nariz, Boca) & & \\ \hline
7. Presentaciones y botellas por lote & & \\ \hline
8. ¿Mostrar ficha completa o como ``Próximamente''? & \uradio\ Ficha completa \quad \uradio\ Próximamente & \uradio\ Ficha completa \quad \uradio\ Próximamente \\
\bottomrule
\end{longtable}

\subsection{Parte D: Datos de Contacto, Ubicación y Cumplimiento Legal}

\begin{pendingbox}{2. Canales de Contacto, Ubicación y Datos Fiscales para Aviso de Privacidad}
\textbf{a) Número oficial de WhatsApp Business (a 10 dígitos):} \texttt{+52} \rule{6cm}{0.5pt}\\[0.5em]
\textbf{b) Correo electrónico oficial donde llegarán los mensajes del formulario:}
\fillline
\textbf{c) Coordenadas o enlace exacto de Google Maps del Palenque en El Gramal:}
\fillline
\textbf{d) Lista de puntos de venta físicos actuales (dirección, nombre del local o ciudad):}
\fillline
\fillline
\textbf{e) Nombre completo del titular o Razón Social (para el Aviso de Privacidad LFPDPPP):}
\fillline
\textbf{f) Domicilio fiscal o comercial para notificaciones de privacidad:}
\fillline
\end{pendingbox}

\subsection{Parte E: Checklist de Entrega de Archivos Multimedia}
\begin{itemize}[label=$\square$, leftmargin=2em, itemsep=0.4em]
    \item \textbf{Logo oficial de Mil Sensaciones} en formato vectorial (\texttt{.SVG}, \texttt{.AI}, \texttt{.EPS}) o \texttt{.PNG} sin fondo.
    \item \textbf{Fotografías en alta resolución de las botellas} (Mezcal Espadín y otras expresiones) bien iluminadas para recorte de fondo transparente.
    \item \textbf{Fotografía real del Maestro Mezcalero Damián Olivera Díaz} en el palenque o campos de agave.
    \item \textbf{Carpeta de fotografías del palenque y proceso artesanal} (corte, cocción, molienda, fermentación, destilación).
    \item \textbf{Archivo de video del proceso de elaboración / palenque} (en formato \texttt{.MP4} / \texttt{.MOV}) para el encabezado principal (Hero).
    \item \textbf{Arte gráfico de las etiquetas de botella} (para referencia de diseño y tipografías).
\end{itemize}

\subsection{Parte F: Decisiones Comerciales Pendientes para la Plataforma de Ventas}

Estas decisiones son necesarias para configurar los módulos RF-16 a RF-27. Sin ellas, el sistema se entrega con valores de ejemplo.

\begin{pendingbox}{3. Precios, Pagos, Envíos y Operación}
\textbf{a) Precio al público por botella de cada expresión (MXN, IVA incluido):}
\fillline
\textbf{b) Precio de mayoreo para bares y restaurantes, y cantidad mínima:}
\fillline
\textbf{c) Formas de pago aceptadas y pasarela elegida:}\\[0.3em]
\uchk\ Tarjeta en línea \quad \uchk\ OXXO \quad \uchk\ SPEI en línea \quad \uchk\ Transferencia manual \quad \uchk\ Efectivo en el palenque\\[0.3em]
Pasarela: \quad \uradio\ Mercado Pago \quad \uradio\ Stripe \quad Correo de la cuenta: \rule{5cm}{0.5pt}
\fillline
\textbf{d) Zonas de envío, paquetería preferida y costo de envío (tarifa fija o cotización por WhatsApp):}
\fillline
\textbf{e) Política de devoluciones o reposición por daño en el transporte:}
\fillline
\textbf{f) Nombres y correos de los Administradores y de los Vendedores POS:}
\fillline
\fillline
\textbf{g) Horario de atención humana para el chat del sitio y WhatsApp:}
\fillline
\textbf{h) Número telefónico dedicado a la API de WhatsApp (no activo en la app):}\\[0.4em]
\texttt{+52} \rule{6cm}{0.5pt}\\[0.6em]
\textbf{i) Datos para la facturación CFDI 4.0:}\\[0.3em]
RFC: \rule{4cm}{0.5pt} \quad Régimen fiscal: \rule{5cm}{0.5pt} \quad C.P. fiscal: \rule{2cm}{0.5pt}\\[0.4em]
\uchk\ CSD vigente (\texttt{.cer}, \texttt{.key} y contraseña) \quad Contador (nombre y correo): \rule{5.5cm}{0.5pt}\\[0.6em]
\textbf{j) Periodo de garantía y esquema de mantenimiento después de la entrega:}
\fillline
\end{pendingbox}

\vspace{1cm}

% ============================================================
% FIRMAS DE CONFORMIDAD
% ============================================================
\noindent
\begin{tabularx}{\linewidth}{@{} X @{\hspace{1.5cm}} X @{}}
    \hrulefill & \hrulefill \\
    \centering\textbf{Por el Cliente: Mezcal Mil Sensaciones}\par
    \small Damián Olivera Díaz / Representante de Marca\par
    \footnotesize El Gramal, Nejapa de Madero, Yautepec, Oaxaca
    &
    \centering\textbf{Elaboró y Especificó (Equipo de Desarrollo)}\par
    \small \textbf{Ing. en Computación Jose Ramon Aragon Toledo}\par
    \footnotesize Ingeniero en Computación — Clave DGP: \texttt{521344}\par
    \footnotesize \texttt{aatr010423@gmail.com}
\end{tabularx}

\end{document}
